% \chapter[Fourier变换和Laplace变换]{傅里叶变换和拉普拉斯变换}
% \chapter[傅里叶变换和拉普拉斯变换]{傅里叶变换和拉普拉斯变换}
% \chapter[积分变换]{积分变换}
\chapter{变分法初步}

\begin{flushright}
\textit{“变则通，通则久。”——《周易·系辞下》}
\footnote{变的是试探路径和试探函数.泛函在一切允许的变分下都取驻值时,那条路径就是所求的极值.}
\end{flushright}

变分法在数学、物理、工程和其它学科有广泛应用价值.这一章我们来介绍泛函和变分法.经典处理见~\cite{gelfand,goldstein,ivrii}.
\section{变分法}


\subsection{泛函}
\label{subsec:functional}
通过以下的几个例子,我们可以了解变分法具体解决什么问题.

如图所示,现在寻找在两点$(x_1, y_1)$和$(x_2,y_2)$间找一个最短的路径$y=u(x)$.常识告诉我们,两点之间直线最短.
数学上表示出来为直线方程,
$$
  y = k x + b = \frac{y_2 - y_1}{x_2 - x_1} ( x - x_1) + y_1  
$$
这样的常识其实在数学上等价为一个最小化问题.对于任意路径,我们利用微积分求得该路径的长度.
$ds^2 = dx^2 + dy^2$,
$$
J[u] = \int_{x_1}^{x_2} \sqrt{ 1 + u'(x)^2} dx    
$$
其中 $u'(x) = du(x)/dx$,满足边界条件$u(x_1) = y_1, u(x_2) = y_2$.因此,我们的问题即求如何找到函数$u(x)$使得
$J$最小化.具体的求解我们下面介绍.这里的方括号区别于函数的圆括号,用来表示这里的映射是函数$\to$数的.
这种$J$和整个函数$u(x)$的依赖关系成为\textbf{泛函关系}.
设对于 (某一函数集合内的) 任意一个函数 $u(x)$, 
有另一个数 $J[u]$ 与之对应, 则称 $J[u]$ 为 $u(x)$ 的\textbf{泛函}. 
这里的函数集合, 即泛函的定义域, 通常包含要求 $u(x)$ 满足一定的边界条件,
并且具有连 续的二阶导数. 这样的 $u(x)$ 称为可取函数. 
% 这种函数的函数称为\textbf{泛函}.

\begin{itemize}
  \item \textbf{最速降线问题}(rachistochrone problem) 一质点在重力作用下无摩擦的沿着某一起终点高度差恒定的轨道下降,初速度为零,如何设计
  轨道形状使得所需下降时间最短.
  \begin{figure}
    \centering
    \includegraphics[width=0.5\linewidth,height=15em]{fastest_track.png}
    \caption{最速降线坐标表示.}
    \label{fig:rachistochrone}
  \end{figure}
  如图\ref{fig:rachistochrone}所示建立坐标,根据动能定理任意时刻的速度可以得到,进而得到总时间
  \begin{equation}
    T[u] = \int_{(x_0,y_0)}^{(x_1, y_1)} \frac{ds}{\sqrt{2 g(y_0 - u)}} = \int_{x_0}^{x_1} \frac{\sqrt{1+u'^2}}{\sqrt{2 g (y_0 - u)}} dx.
  \end{equation}

  \item \textbf{最小作用量原理}(least action principle) 理论力学上,它可以导出牛顿力学方程.我们定义的体系的动能为$K = \half m \bv^2$,
  势能为$V = V(\br)$, 定义作用量
  \begin{equation}
    A[\br] = \int \left[ \half m \left( \frac{d\br}{dt} \right)^2 - V(\br) \right] dt 
  \end{equation}
  上式利用了定义$\vec{\bv} = \frac{d\br}{dt}$. 对作用量$A$进行最小化将得到$\br$满足的方程,即牛顿第二定律方程.

  \item \textbf{悬链线问题} (catenary problem) \quad 自然界中有许多出现悬链线形状的地方,如雨后的蜘蛛网, 路边的铁链, 高压输电线等(如图\ref{fig:catenary_example}), 有一质地均匀、柔软的绳索,两端固定,绳索仅受重力的作用而下垂.试问该绳索在平衡状态时是怎样的曲线?
  % 设绳索的最低点为$A$ .设绳索曲线的方程为$y=\phi(x)$,
  % 考虑绳索上点$A$到另一点$M$间的一段弧$AM$,设其长为s.假定绳索的线密度为$\rho$,则弧$AM$所受重力为$\rho gs$.由于绳索是柔软的,因而在点$A$
  % 处的张力沿水平的切线方向(该点导数为零),设其大小设为$H$;在点$M$处的张力沿该点处的切线方向,设其倾角为$\theta$,
  % 其大小为$T$,因为作用于弧$AM$的外力相互平衡,把作用于弧$AM$上的力沿竖直及水平两方向分解,得:
  % $$T \sin \theta=\rho g s, T \cos \theta=H$$
  % 于是有,
  % $$
  %   y^{\prime}=a \int_0^x \sqrt{1+y^{\prime 2}} d x
  % $$ 
  \begin{figure}
    \centering

    \includegraphics[width=0.45\linewidth]{roadside.png}
    \includegraphics[width=0.45\linewidth]{cobweb.png}
    \caption{雨后蜘蛛网和路边的铁链.}
    \label{fig:catenary_example}
  \end{figure}

  假设重力加速度恒为 $\mathrm{g}$ ,则某物体重力势能可表述为 $U=m g h$. 设绳子两端点距离为 $\mathrm{L}$.
建立 $x-y$ 坐标系, $y$ 表示高度, $y(x)$ 代表绳索形成的曲线,绳索关于 $y$ 轴对称.绳子上每个点的重力势 能个表示为重力势能公式两边求导的结果： $d U=d(m g \cdot y(x))=g \cdot y(x) \cdot d m$.
设绳索单位长度质量为 $\rho$ ,则根据其定义, $\frac{d m}{d s}=\rho$ ,所以
$d m=\rho \sqrt{d x^2+d y^2}=\rho \sqrt{1+\left(\frac{d y}{d x}\right)^2} d x$
代入得 $d U=\rho g \sqrt{1+\left(\frac{d y}{d x}\right)^2} d x \cdot y(x)$
作用于整条绳索上的重力势能表示为:
$$
U[y(x)]=\int_{-L / 2}^{L / 2} \rho g \sqrt{1+y^{\prime}(x)^2} \cdot y(x) d x,
$$
悬链线就是在此情形中重力势能最小的曲线. 


\end{itemize}

% \subsection{泛函导数}
% 对于一般的情况


\subsection{泛函的极值}

对于一般的情况,有函数到数的映射写成
\begin{equation}
  J[u] = \int_{a}^{b}  dx L(u, u', u'', \cdots| x), 
\end{equation}
其中$L(u, u', u'',\cdots | x)$为$u$的半局域泛函, $u' \equiv du/dx, u'' \equiv = du^2/dx^2$.
对于不包含一阶以上的导数时,即$L = L(u,u')$,对函数$u(x)$作一微小变化$\delta u(x)$,称为$u(x)$的\textbf{变分},在积分上下限处,
有$\delta u(a) = \delta u(b) = 0$. $J[u]$取最小值时,需满足$\delta J[u] \equiv = J[u + \delta u] - J[u]$
为零.利用$\delta u' = \funcpd{}{x} \delta u$和$fg' = (fg)' - f'g$,
\begin{align}
  J[u + \delta u] &= \int_{a}^{b}  dx L(u + \delta u, u' + \delta u'| x) 
    \nonumber \\ 
    &= \int_{a}^{b} dx \left[ L(u, u'| x) + \funcpd{L}{u} \delta u + \funcpd{L}{u'} \delta u' \right]  
    \nonumber \\ 
    &= \int_{a}^{b} dx \left[ L(u, u'| x) + \funcpd{L}{u} \delta u + \funcpd{}{x}\left( \funcpd{L}{u'} \delta u \right)
    -  \funcpd{}{x}\left( \funcpd{L}{u'} \right) \delta u  \right]  
    \nonumber \\ 
    &= \int_{a}^{b} dx \left[ L(u, u'| x) + \left( \funcpd{L}{u} -  \funcpd{}{x}\left( \funcpd{L}{u'} \right)   \right) \delta u \right]  
        + \left. \funcpd{L}{u'} \delta u \right|_{a}^{b} 
     \nonumber %\\ 
\end{align}
于是,结合边界条件我们有
\begin{align}
  \delta J[u] = \int_{a}^{b} dx \left[ \funcpd{L}{u} -  \funcpd{}{x}\left( \funcpd{L}{u'} \right)   \right] \delta u 
\end{align}
被积函数对任意变化$\delta u$都需要满足$\delta J[u] = 0$,于是我们得到
\begin{equation}
  \funcpd{L}{u} -  \funcpd{}{x}\left( \funcpd{L}{u'} \right) = 0 .
\end{equation}
这就是著名的\textbf{欧拉-拉格朗日方程}(Euler-Lagrange Equation).

利用该方程,可以得到两点之间距离最短满足
$$
  0 = - \funcd{}{x} \frac{u'}{\sqrt{1 + u'^2}} = \frac{-u''}{(1+ u'^2)^{3/2}}
$$
也就是
$$ u'' = 0$$
也就是解为直线方程$ u = k x + b$,$k,b$的值由边界条件确定.

上面从 $\delta J=0$ 推出欧拉--拉格朗日方程,用到的是\textbf{变分法基本引理}:
若 $M(x)$ 连续,并且对一切在端点为零的光滑变分 $\delta u$ 都有
\begin{equation}
  \int_a^b M(x)\,\delta u(x)\,\mathrm{d}x = 0,
\end{equation}
则 $M(x)=0$. 泛函取极值只保证积分等于零;引理把这个积分条件变成逐点的微分方程.

若 $L$ 不显含自变量 $x$,欧拉--拉格朗日方程有首次积分
\begin{equation}
  L - u' \funcpd{L}{u'} = C .
  \label{eq:beltrami}
\end{equation}
最短路径的 $L=\sqrt{1+u'^2}$ 就不含 $u$,于是 $u'/\sqrt{1+u'^2}$ 为常数,再求导仍得到 $u''=0$.
下面的最速降线和悬链线都靠这个首次积分,而不必把欧拉--拉格朗日方程展开成二阶方程.

\begin{example}
最速降线的解是旋轮线.
\end{example}
\begin{solution}
质点从高度 $y_0$ 由静止下滑.令下落距离 $v(x)=y_0-u(x)\ge 0$,则 $u'=-v'$,时间泛函写成
\begin{equation}
  T[v] = \int_{x_0}^{x_1} \frac{\sqrt{1+v'^2}}{\sqrt{2gv}}\,\mathrm{d}x .
\end{equation}
被积函数不显含 $x$.由\eqref{eq:beltrami},
\begin{equation}
  \frac{1}{\sqrt{1+v'^2}\,\sqrt{2gv}} = C ,
\end{equation}
即 $v'^2 = (2a-v)/v$,其中正的常数 $2a=1/(2gC^2)$.
用参数 $v=a(1-\cos\theta)$,$\mathrm{d}v=a\sin\theta\,\mathrm{d}\theta$.
此时 $\sqrt{v/(2a-v)}=\sin(\theta/2)/\cos(\theta/2)$,而 $\sin\theta=2\sin(\theta/2)\cos(\theta/2)$,所以
\begin{equation}
  \mathrm{d}x = \sqrt{\frac{v}{2a-v}}\,\mathrm{d}v = a(1-\cos\theta)\,\mathrm{d}\theta .
\end{equation}
积分得
\begin{equation}
  x-x_0 = a(\theta-\sin\theta), \qquad v = a(1-\cos\theta).
\end{equation}
这是旋轮线,不是直线.直线使路程最短,旋轮线使时间最短;泛函不同,极值曲线就不同.
常数 $a$ 由终点决定.
\end{solution}

取终点对应参数从 $0$ 到 $\pi$,两端为 $(x_0,y_0)$ 与 $(x_0+\pi R,\, y_0-2R)$,这里 $R$ 即上面的 $a$.弧长与速度是
\begin{equation}
  \mathrm{d}s = 2R\sin\frac{\varphi}{2}\,\mathrm{d}\varphi,
  \qquad
  v = 2\sin\frac{\varphi}{2}\,\sqrt{gR},
\end{equation}
所以摆线用时是常数被积函数的积分,
\begin{equation}
  T_{\textrm{摆线}} = \int_0^{\pi}\sqrt{\frac{R}{g}}\,\mathrm{d}\varphi = \pi\sqrt{\frac{R}{g}}.
\end{equation}
连接同一对端点的直线,水平距离 $\pi R$、落差 $2R$,斜率 $-2/\pi$.沿直线积分得到
\begin{equation}
  T_{\textrm{直线}} = \sqrt{\pi^2+4}\,\sqrt{\frac{R}{g}}.
\end{equation}
二者之比 $\sqrt{\pi^2+4}/\pi\approx 1.185$.摆线用时约为直线的 $84\%$.

\subsection{高阶导数与自然边界条件}
若拉格朗日量含到 $m$ 阶导数,
\begin{equation}
  J[u] = \int_a^b L\bigl(x,u,u',\ldots,u^{(m)}\bigr)\,\mathrm{d}x,
\end{equation}
一阶变分是
\begin{equation}
  \delta J = \int_a^b \sum_{k=0}^{m} \funcpd{L}{u^{(k)}}\,\delta u^{(k)}\,\mathrm{d}x.
\end{equation}
变分与求导可交换,$\delta u^{(k)}=\mathrm{d}^k(\delta u)/\mathrm{d}x^k$.对含 $\delta u^{(k)}$ 的项作 $k$ 次分部积分,内部项变成
\begin{equation}
  (-1)^k \int_a^b \frac{\mathrm{d}^k}{\mathrm{d}x^k}\left(\funcpd{L}{u^{(k)}}\right)\delta u\,\mathrm{d}x.
\end{equation}
对任意内部变分 $\delta J=0$,就得到高阶欧拉--拉格朗日方程
\begin{equation}
  \fnald{J}{u}
  = \sum_{k=0}^{m}(-1)^k \frac{\mathrm{d}^k}{\mathrm{d}x^k}\left(\funcpd{L}{u^{(k)}}\right)
  = 0.
\end{equation}
$m=1$ 时就是前面的标准形式.

分部积分留下的边界项是 $\delta u,\delta u',\ldots,\delta u^{(m-1)}$ 的线性组合.端点上有三种典型设定.
\begin{itemize}
  \item 狄利克雷:$u(a),u(b)$ 给定,则 $\delta u(a)=\delta u(b)=0$,一阶情形的边界项自动消失.
  \item 诺伊曼:端点的 $u$ 自由,则必须另有
  \begin{equation}
    \left.\funcpd{L}{u'}\right|_{a}
    = \left.\funcpd{L}{u'}\right|_{b} = 0,
  \end{equation}
  边界项才为零.
  \item 混合:有的端点固定函数值,有的端点固定导数,条件按端点分别写.
\end{itemize}
$m>1$ 时,每一个自由的高阶导数 $\delta u^{(j)}$ 都再给出一个条件.这些不是另外假定的,而是 $\delta J=0$ 对自由边界变分的要求,称为\textbf{自然边界条件}.

多个因变量 $u_1,\ldots,u_n$ 的变分彼此独立,每个 $u_i$ 各有一个欧拉--拉格朗日方程.场 $u(\bx)$ 的泛函 $J[u]=\int_{\Omega} F(\bx,u,\partial_i u)\,\mathrm{d}^d\bx$ 同样给出
\begin{equation}
  \funcpd{F}{u} - \partial_i\left(\funcpd{F}{\partial_i u}\right) = 0,
\end{equation}
只是把 $\mathrm{d}/\mathrm{d}x$ 换成了散度.

对最小作用量原理问题,可以得到
$$
  0 = - \funcd {V(\br)}{\br} - \funcd{}{t} \left[ m \frac{d\br}{dt} \right]
$$
若质量$m$不依赖于时间,得到牛顿方程,
$$
  m \frac{d^2\br}{dt^2} = - \funcd {V(\br)}{\br} \equiv \bF .
$$

对于如多目标变量的$J(u_1, u_2, \cdots, u_n, u_1', u_2', \cdots, u_n'|x)$, 对于每一个$u_i$, $i=1,2,\dots,n$都有独立的欧拉方程

\begin{equation}
  \funcpd{L}{u_i} -  \funcpd{}{x}\left( \funcpd{L}{u'_i} \right) = 0 .
\end{equation}

对于多变量的情况,$u(x,y)$, 如
\begin{equation}
J[u]=\iint_R F\left(x, y, u, u_x, u_y\right) d x d y .
\end{equation}
对应的Euler方程为
\begin{equation}
F_u-\frac{\partial}{\partial x} F_{u_x}-\frac{\partial}{\partial y} F_{u_y}=0.
\end{equation}
例如,对于肥皂泡构成的膜,物理上表面张力要求该膜形成的膜面积最小, 对应的能量可以用$E(u)$来表示,
用合适的单位,可以表示
$$
E(u)=\iint_S\left[1+\left(\frac{\partial u}{\partial x}\right)^2+\left(\frac{\partial u}{\partial y}\right)^2\right]^{1 / 2} d x d y.
$$
Euler方程为最小面积方程
$$
-\frac{\partial}{\partial x}\left(\frac{1}{\sqrt{A}} \frac{\partial u}{\partial x}\right)-\frac{\partial}{\partial y}\left(\frac{1}{\sqrt{A}} \frac{\partial u}{\partial y}\right)=0,
$$
其中
$$
A=1+\left(\frac{\partial u}{\partial x}\right)^2+\left(\frac{\partial u}{\partial y}\right)^2.
$$
它表示平衡曲面的平均曲率为零.未知量现在依赖两个自变量,所以变分条件是偏微分方程,不再是常微分方程.

\subsection{泛函导数}
把变分写成 $u_\epsilon=u+\epsilon\eta$,泛函导数由
\begin{equation}
  \left.\funcd{}{\epsilon} J[u+\epsilon\eta]\right|_{\epsilon=0}
  = \int_a^b \fnald{J}{u(x)}\,\eta(x)\,\mathrm{d}x
\end{equation}
定义.固定端点时,一阶导数的拉格朗日量给出
\begin{equation}
  \fnald{J}{u}
  = \funcpd{L}{u} - \funcd{}{x}\left(\funcpd{L}{u'}\right).
\end{equation}
欧拉--拉格朗日方程就是这个泛函导数为零.

非局域泛函把不同点的函数值耦在一起.例如
\begin{equation}
  J[u] = \half \iint u(x)\,K(x,x')\,u(x')\,\mathrm{d}x\,\mathrm{d}x'.
\end{equation}
$K$ 对称时 $\fnald{J}{u(x)}=\int K(x,x')u(x')\,\mathrm{d}x'$;$K$ 不对称时核要换成 $\half\bigl(K(x,x')+K(x',x)\bigr)$.下面的哈特里项就是这种结构.

若能量是密度 $n(\br)$ 的泛函,并且粒子数 $\int n(\br)\,\mathrm{d}^3\br=N$ 固定,引入化学势 $\mu$,
\begin{equation}
  \widetilde E[n] = E[n] - \mu\left(\int n\,\mathrm{d}^3\br - N\right).
\end{equation}
对 $n$ 变分得到
\begin{equation}
  \fnald{E}{n(\br)} = \mu .
\end{equation}
$\mu$ 在空间中是常数,尽管外势和密度都可以随地点变化.

\begin{example}
局域密度泛函 $E[n]=\int\bigl(C n^{5/3}+v(\br)n\bigr)\,\mathrm{d}^3\br$ 的泛函导数.
\end{example}
\begin{solution}
被积函数不含 $n$ 的导数,所以
\begin{equation}
  \fnald{E}{n(\br)} = \frac{5}{3} C n(\br)^{2/3} + v(\br).
\end{equation}
归一化后的极值条件是 $\tfrac{5}{3}C n^{2/3}+v=\mu$,在每一点都是代数方程,只通过 $\mu$ 和总粒子数发生整体联系.
\end{solution}

若再加入电子之间的库仑能
\begin{equation}
  E_H[n] = \half \iint \frac{n(\br)n(\brp)}{|\br-\brp|}\,\mathrm{d}^3\br\,\mathrm{d}^3\brp,
\end{equation}
则
\begin{equation}
  \fnald{E_H}{n(\br)} = \int \frac{n(\brp)}{|\br-\brp|}\,\mathrm{d}^3\brp .
\end{equation}
局域动能项和外势在同一点取值,哈特里项则把整个空间的密度都记进来.把三项放在一起,
\begin{equation}
  E[n] = C\int n^{5/3}\,\mathrm{d}^3\br
  + \int v_{\mathrm{ext}}(\br)\,n\,\mathrm{d}^3\br
  + E_H[n],
\end{equation}
其中 Thomas--Fermi 系数 $C=\frac{3}{10}(3\pi^2)^{2/3}$.极值条件是
\begin{equation}
  \frac{5}{3} C\, n(\br)^{2/3} + v_{\mathrm{ext}}(\br) + v_H(\br) = \mu,
  \qquad
  v_H(\br) = \int \frac{n(\brp)}{|\br-\brp|}\,\mathrm{d}^3\brp.
\end{equation}
粒子数守恒用来定 $\mu$.








\subsection{哈密顿方程}
\idx{哈密顿方程}{hamidunfangcheng}
经典物理里,拉式量$L(x_1, x_2, \cdots, x_n, \dot{x}_1, \dot{x}_2, \cdots, \dot{x}_n)|t)$是位置和速度的函数.\textbf{哈密顿原理}
说的是系统的运动从$t_1$到$t_2$使得拉式量的时间积分的变分为零.即
\begin{equation}
    \delta \int_{t_1}^{t_2} L\left(x_1, x_2, \ldots, x_n, \dot{x}_1, \dot{x}_2, \ldots, \dot{x}_n | t\right) d t=0 .
\end{equation}
于是有拉格朗日运动方程

\begin{equation}
    \frac{d}{d t} \frac{\partial L}{\partial \dot{x}_i}-\frac{\partial L}{\partial x_i}=0, \quad \forall i
\end{equation}

拉格朗日运动方程同牛顿运动方程是等价的.拉格朗日运动方程有一些优点.首先,依赖坐标不必是标准的坐标或长度,
它们可以根据物理问题具体选择,拉式运动方程对坐标系统是不变的.其次,拉式运动方程可以扩展到各个物理分支,如电磁场
和量子力学以及更远的规范场论等.此外,对于拉式量的变换,守恒律和对称性可以很容易通过诺特定理对应起来.


哈密顿注意到拉式运动方程可以简化为一系列偏微分方程,称为\textbf{哈密顿方程}.
正则动量定义为
\begin{equation}
    p_i=\frac{\partial L}{\partial \dot{q}_i}
\end{equation}

\begin{equation}
    d L=\sum_i\left(\frac{\partial L}{\partial q_i} d q_i+\frac{\partial L}{\partial \dot{q}_i} d \dot{q}_i\right)+\frac{\partial L}{\partial t} d t=\sum_i\left(\dot{p}_i d q_i+p_i d \dot{q}_i\right)+\frac{\partial L}{\partial t} d t
\end{equation}

定义体系哈密顿量为
\begin{equation}
    H(q_i, p_i, t)=\sum_i p_i \dot{q}_i-L 
\end{equation}
我们有

\begin{equation}
    d H=\sum_i\left(p_i d \dot{q}_i+\dot{q}_i d p_i\right)-\left(\sum_i\left(\dot{p}_i d q_i+p_i d \dot{q}_i\right)+\frac{\partial L}{\partial t} d t\right)=\sum_i\left(\dot{q}_i d p_i-\dot{p}_i d q_i\right)-\frac{\partial L}{\partial t} d t
\end{equation}

另外根据定义有
\begin{equation}
    d H=\sum_i\left(\frac{\partial H}{\partial p_i} d p_i+\frac{\partial H}{\partial q_i} d q_i\right)+\frac{\partial H}{\partial t} d t .
\end{equation}
于是得到哈密顿方程

\begin{align}
    \frac{\partial H}{\partial p_i}&=\dot{q}_i
    \\
    \frac{\partial H}{\partial q_i}&=-\dot{p}_i
    \\
    \frac{\partial H}{\partial t}&=-\frac{\partial L}{\partial t}
\end{align}
对于保守系统,$\frac{\partial H}{\partial t}=0$, $H$为守恒量对应体系的能量.

沿真实轨道 $q_i(t),p_i(t)$ 对 $H(q_i,p_i,t)$ 求全导数,
\begin{equation}
  \funcd{H}{t}
  = \sum_i\left( \funcpd{H}{q_i}\dot q_i + \funcpd{H}{p_i}\dot p_i \right)
    + \funcpd{H}{t}.
\end{equation}
代入哈密顿方程 $\dot q_i=\partial H/\partial p_i$、$\dot p_i=-\partial H/\partial q_i$,前两项相消,
\begin{equation}
  \funcd{H}{t} = \funcpd{H}{t} = -\funcpd{L}{t}.
\end{equation}
因此只要 $L$ 不显含时间,$\mathrm{d}H/\mathrm{d}t=0$,$H$ 沿轨道是常数.这对应时间平移对称性.对通常的 $L=K-V$,这个常数就是总能量.动量、角动量等其他守恒量同样来自 $L$ 的对称性.

\begin{example}
自由粒子与简谐振子的哈密顿量.
\end{example}
\begin{solution}
自由粒子 $L=\half m\dot q^2$.共轭动量 $p=m\dot q$,所以 $H=p^2/(2m)$.哈密顿方程是 $\dot q=p/m$、$\dot p=0$,即匀速直线运动.

简谐振子 $L=\half m\dot q^2-\half k q^2$.仍有 $p=m\dot q$,而
\begin{equation}
  H = \frac{p^2}{2m} + \half k q^2.
\end{equation}
哈密顿方程 $\dot q=p/m$、$\dot p=-kq$ 合起来就是 $m\ddot q+kq=0$.$L$ 不含 $t$,所以 $H$ 守恒,等于动能加势能.
\end{solution}
\subsection{拉格朗日乘子}
\idx{拉格朗日乘子}{lagelangrichengzi}
积分约束是把整条曲线限制住.更基本的情形是函数 $f(x_1,\ldots,x_n)$ 在若干等式 $c_a(x_1,\ldots,x_n)=0$ 下取极值.引入乘子 $\lambda_a$,作
\begin{equation}
  \mathcal{F} = f - \sum_a \lambda_a c_a .
\end{equation}
把 $x_i$ 和 $\lambda_a$ 都当成独立变量,驻点方程是
\begin{equation}
  \funcpd{f}{x_i} - \sum_a \lambda_a \funcpd{c_a}{x_i} = 0,
  \qquad c_a = 0 .
\end{equation}
对 $\lambda_a$ 求导只是把约束写回来.

几何上,$\nabla f$ 垂直于 $f$ 的等值面.约束曲面上允许的位移 $\mathrm{d}\bx$ 满足 $\nabla c_a\cdot\mathrm{d}\bx=0$.
要让 $f$ 沿这些位移的一阶变化为零,$\nabla f$ 必须落在各个 $\nabla c_a$ 张成的法向里:
\begin{equation}
  \nabla f = \sum_a \lambda_a \nabla c_a .
\end{equation}
沿约束面走动时,$f$ 的一阶变化为零;在整个空间里 $\nabla f$ 不必为零.

\begin{example}
求平面 $x+y+z=1$ 上离原点最近的点.
\end{example}
\begin{solution}
极小化 $f=x^2+y^2+z^2$,约束 $c=x+y+z-1=0$.驻点方程为
\begin{equation}
  2x=\lambda,\qquad 2y=\lambda,\qquad 2z=\lambda,\qquad x+y+z=1.
\end{equation}
于是 $x=y=z$,最近点是 $(1/3,1/3,1/3)$.三个偏导相等,是因为平面的法向是 $(1,1,1)$,而 $\nabla f=(2x,2y,2z)$ 必须与它平行.
\end{solution}

\begin{example}
求 $f=x+y$ 在双曲线 $xy=1$ 上的驻点.
\end{example}
\begin{solution}
\begin{equation}
  1-\lambda y = 0, \qquad 1-\lambda x = 0, \qquad xy=1.
\end{equation}
前两式给出 $x=y$.代入 $xy=1$,得到 $(1,1)$ 和 $(-1,-1)$.
在这两点,$x+y$ 的等值线与双曲线相切.
\end{solution}

\begin{example}
单位球面上 $f=z$ 的极值.
\end{example}
\begin{solution}
约束 $c=x^2+y^2+z^2-1=0$,增广函数 $\mathcal{F}=z-\lambda c$.
\begin{equation}
  -2\lambda x=0,\qquad -2\lambda y=0,\qquad 1-2\lambda z=0,\qquad x^2+y^2+z^2=1.
\end{equation}
若 $\lambda=0$,第三式变成 $1=0$,矛盾.因此 $x=y=0$,再由球面得 $z=\pm 1$,$\lambda=1/(2z)$.极值点是北极 $(0,0,1)$ 和南极 $(0,0,-1)$.几何上 $\nabla f=(0,0,1)$ 必须平行于 $\nabla c=(2x,2y,2z)$,球面上只有这两点满足.
\end{solution}

\begin{example}
单位圆上内接 $N$ 边形的最大面积.
\end{example}
\begin{solution}
顶点 $(x_i,y_i)$ 在单位圆上,$x_i^2+y_i^2=1$,按逆时针编号,并令 $(x_{N+1},y_{N+1})=(x_1,y_1)$.由格林公式,围成的有向面积是
\begin{equation}
  2A = \sum_{i=1}^{N}\bigl(x_i y_{i+1}-x_{i+1} y_i\bigr).
\end{equation}
写成 $x_i=\cos\theta_i$,$y_i=\sin\theta_i$,每一项是相邻圆心角的正弦,
\begin{equation}
  A = \half\sum_{i=1}^{N}\sin(\theta_{i+1}-\theta_i).
\end{equation}
对 $\theta_k$ 求导,它只出现在第 $k-1$ 项和第 $k$ 项里,
\begin{equation}
  \funcpd{A}{\theta_k}
  = \half\Bigl[\cos(\theta_k-\theta_{k-1})-\cos(\theta_{k+1}-\theta_k)\Bigr]=0.
\end{equation}
记 $\alpha_i=\theta_{i+1}-\theta_i\in(0,\pi)$,则 $\cos\alpha_{k-1}=\cos\alpha_k$,从而各圆心角相等.它们的和是 $2\pi$,所以 $\alpha_i=2\pi/N$.面积最大的是正 $N$ 边形,
\begin{equation}
  A_{\max} = \frac{N}{2}\sin\frac{2\pi}{N}.
\end{equation}
$N\to\infty$ 时 $A_{\max}\to\pi$,即单位圆的面积.

用 $z_i=x_i+\mathrm{i} y_i$,在单位圆上 $\bar z_i=1/z_i$,同一面积可写成
\begin{equation}
  2A = \operatorname{Im}\sum_{i=1}^{N}\bar z_i z_{i+1}.
\end{equation}
$z_i=e^{\mathrm{i}\theta_i}$ 时 $\bar z_i z_{i+1}=e^{\mathrm{i}(\theta_{i+1}-\theta_i)}$,虚部就是上面的正弦.
\end{solution}

\subsection{约束问题}
我们常常遇到的一类问题是$u$满足其他的约束条件,比如上述两点之间曲线长度问题可以增加一个约束条件,
所构成的曲线所围面积为定值$A$,即$\int u(x) dx = A$.通常的做法是利用拉格朗日乘子法,即对
目标函数$J[u]$添加一项$\lambda\left( \int u(x) dx - A \right)$,变成无约束问题去求解. 
这里的$\lambda$称为\textbf{拉格朗日乘子}.
以上例为例,我们需要极小化的函数变成
$$
J[u] = \int_{x_1}^{x_2} \sqrt{ 1 + u'(x)^2} dx   + \lambda \left( \int_{x_1}^{x_2} u(x) dx - A \right).
$$
这样拉式量内包含了这一约束.对于该例子应用欧拉-拉格朗日方程,我们有
\[
  \lambda - \funcd{}{x} \frac{u'}{\sqrt{1 + u'^2}} = 0,
\]
可以求得
\[
 \frac{u'}{\sqrt{1+u'^2}}= \lambda x + c   %\textrm{或} u'=\frac{\lambda x-c}{\sqrt{1-(\lambda x-c)^2}}
 \]
令辅助函数$ p(x)\quad  = \lambda x + c   $,  两边平方,得到
\[ u' = \frac{p}{\sqrt{1 - p^2}} .
\]
换元令$ t = \lambda x + c$, 则有 $dt = \lambda dx$, 于是
\[
 du = u' dx = \frac{t}{\sqrt{1 - t^2}} \frac{dt}{\lambda} .
\]
积分得到
\[
u(x) = \int \frac{t}{\sqrt{1 - t^2}} \frac{dt}{\lambda} = - \frac{\sqrt{1 - t^2}}{\lambda} + d,
\]
其中$d$为积分常数.
最终得到方程
\[
  (\lambda x +c)^2 + (\lambda u -d)^2 = 1
\]
其中$c,d,\lambda$由$u(x_1)=y_1, u(x_2) = y_2$ 和$\int u(x) dx = A$决定.
在这样的约束条件下,最短路径是一段圆弧！
此外, 可以对$p(x)$直接求导得到
\[
\lambda = \frac{u''}{(1+y'^2) ^{3/2}}
\]
而右边正式平面曲线的曲率$\kappa$, 因此在约束条件下, 最短路径的曲率是常数.
故也可知曲线为圆弧.

\begin{example}
悬链线.
\end{example}
\begin{solution}
均匀软绳的重力势能是
\begin{equation}
  U[u] = \int_{-L/2}^{L/2} \rho g\, u\sqrt{1+u'^2}\,\mathrm{d}x,
\end{equation}
绳长另有约束 $\int_{-L/2}^{L/2}\sqrt{1+u'^2}\,\mathrm{d}x=\ell$.
乘子 $\lambda$ 把被积函数改成 $(\rho g u+\lambda)\sqrt{1+u'^2}$.
它不显含 $x$,首次积分给出
\begin{equation}
  \frac{\rho g u+\lambda}{\sqrt{1+u'^2}} = A .
\end{equation}
令 $w=\rho g u+\lambda$,则 $\mathrm{d}w/\mathrm{d}x = (\rho g/A)\sqrt{w^2-A^2}$.
分离变量并积分,
\begin{equation}
  w = A \cosh\Bigl(\frac{\rho g}{A}(x-x_c)\Bigr).
\end{equation}
因此
\begin{equation}
  u(x) = a\cosh\Bigl(\frac{x-x_c}{a}\Bigr) + u_c,
  \qquad a=\frac{A}{\rho g}.
\end{equation}
平衡曲线是双曲余弦,不是抛物线.顶点附近把余弦展开到二阶,才近似成抛物线.常数由两端位置和绳长决定.
\end{solution}


% \subsection{Euler-Lagrange方程}
% 上式的求解其实可以归结

% \subsubsection{二维问题}

\subsection{自由端点}
到这里为止,端点上的函数值都事先给定,变分在端点为零.若右端点的值也参加变分,作用量还可以带一项只依赖端点值的 $M(u(b))$,
\begin{equation}
  J[u]=\int_a^b L(x,u,u')\,\mathrm{d}x+M\bigl(u(b)\bigr),
\end{equation}
左端仍然固定,$\delta u(a)=0$.分部积分后多出来的边界项是
\begin{equation}
  \Biggl(\funcpd{L}{u'}+\frac{\mathrm{d}M}{\mathrm{d}u}\Biggr)_{x=b}\delta u(b).
\end{equation}
$\delta u(b)$ 任意,所以除了欧拉--拉格朗日方程,还要
\begin{equation}
  \left.\funcpd{L}{u'}\right|_{b}+\frac{\mathrm{d}M}{\mathrm{d}u}\bigl(u(b)\bigr)=0.
  \label{eq:free-end}
\end{equation}
左端自由时符号相反.两端都自由,就有两个这样的条件.每个端点恰好一条信息:要么给定 $u$,要么给定\eqref{eq:free-end}.

\begin{example}
牛顿最小阻力问题.旋转体的长度是 $a$,底半径是 $H$,母线 $y(x)$ 从 $y(0)=H$ 画到 $x=a$.在合适的单位下,侧面阻力和前端圆盘的阻力合成
\begin{equation}
  S[y]=\int_0^a \frac{y\,(y')^3}{1+(y')^2}\,\mathrm{d}x+\frac12 y(a)^2.
\end{equation}
$y'$ 取负,$y$ 向右变细.
\end{example}
\begin{solution}
被积函数不显含 $x$.令 $p=-y'>0$,贝尔特拉米首次积分给出
\begin{equation}
  y=C\frac{(1+p^2)^2}{p^3}.
\end{equation}
把 $y$ 看成 $p$ 的函数,
\begin{equation}
  \mathrm{d}x=-\frac{\mathrm{d}y}{p}
  =C\Biggl(\frac{3}{p^5}+\frac{2}{p^3}-\frac{1}{p}\Biggr)\mathrm{d}p,
\end{equation}
积分后 $x$ 和 $y$ 都由 $p$ 参数表示.经典解的前端不是尖的,而是半径为 $y(a)$ 的圆盘,圆盘边缘处母线斜率取 $-1$,所以 $p$ 从 $x=0$ 处的某个 $p_0\in(0,1)$ 增加到 $1$.这个 $-1$ 不是光滑自由端条件\eqref{eq:free-end}:端点项 $M=\half y^2$ 要求 $\partial L/\partial y'+y=0$,而
\begin{equation}
  \funcpd{L}{y'}=y\frac{(y')^2\bigl(3+(y')^2\bigr)}{\bigl(1+(y')^2\bigr)^2}
\end{equation}
对一切实斜率都和 $y$ 同号,两者之和不能是零.母线因此不能以有限斜率光滑地自由终止,圆盘是角点,不是\eqref{eq:free-end}的解.
\end{solution}

多个未知函数 $u_1,\ldots,u_m$ 各自有一条欧拉--拉格朗日方程.若 $L$ 不显含某个 $u_k$,对应的 $\partial L/\partial u_k'$ 是常数.若 $L$ 不显含自变量,则
\begin{equation}
  \sum_{k=1}^{m}u_k'\funcpd{L}{u_k'}-L=C.
\end{equation}
复值函数看成实部和虚部两个实函数,方程数目加倍.

带矢量势的拉格朗日量 $L=K+\mathbf{A}(\mathbf{q})\cdot\dot{\mathbf{q}}-V$ 比 $K-V$ 多一项.欧拉--拉格朗日方程里多出洛伦兹力,力的分量是 $F_{jk}\dot q_k$,
\begin{equation}
  F_{jk}=\partial_j A_k-\partial_k A_j.
\end{equation}
三维时 $\mathbf{F}=\nabla\times\mathbf{A}$,这项就是 $\dot{\mathbf{q}}\times\mathbf{B}$.$\mathbf{A}$ 加上任意函数的梯度不改变方程,只把动量移动一个梯度,作用量改变端点上的值.

\subsection{二维自然边界}
区域 $\Omega$ 上
\begin{equation}
  J[u]=\iint_{\Omega}F(x,y,u,u_x,u_y)\,\mathrm{d}x\,\mathrm{d}y.
\end{equation}
内部变分给出
\begin{equation}
  F_u-\partial_x F_{u_x}-\partial_y F_{u_y}=0.
\end{equation}
边界上若 $u$ 没有给定,分部积分留下
\begin{equation}
  \int_{\partial\Omega}\bigl(F_{u_x}\nu_x+F_{u_y}\nu_y\bigr)\delta u\,\mathrm{d}s=0,
\end{equation}
$\boldsymbol{\nu}$ 是外法向.自然边界条件是 $F_{u_x}\nu_x+F_{u_y}\nu_y=0$.只在边界的一部分上给定 $u$ 时,自然条件写在剩下的那一部分.

\begin{example}
狄利克雷积分
\begin{equation}
  J[u]=\iint_{\Omega}\Bigl(\half|\nabla u|^2-fu\Bigr)\,\mathrm{d}x\,\mathrm{d}y.
\end{equation}
\end{example}
\begin{solution}
$F_{u_x}=u_x$, $F_{u_y}=u_y$, $F_u=-f$.欧拉--拉格朗日方程是 $-\Delta u=f$.边界项是 $\int(\partial u/\partial\nu)\delta u\,\mathrm{d}s$.边界上 $u$ 给定时这一项是零;边界自由时,自然条件是 $\partial u/\partial\nu=0$.泊松方程的诺伊曼条件就是这个自由边界变分,不是另外加进去的.
\end{solution}

极小曲面 $z=u(x,y)$ 的面积是 $\iint\sqrt{1+u_x^2+u_y^2}\,\mathrm{d}x\,\mathrm{d}y$,欧拉--拉格朗日方程就是平均曲率为零,前面已经写过.若再减去外力做的功 $\iint f u\,\mathrm{d}x\,\mathrm{d}y$,方程右边多出 $f$,边界上 $u$ 仍按给定的边线固定.

旋转曲面 $z=u(r)$ 的面积化成单积分
\begin{equation}
  S[u]=2\pi\int r\sqrt{1+u_r^2}\,\mathrm{d}r.
\end{equation}
被积函数不含 $u$,所以 $r u_r/\sqrt{1+u_r^2}=C$,
\begin{equation}
  u(r)=C\,\mathrm{acosh}\frac{r}{C}+D.
\end{equation}
极小旋转曲面是悬链面.两端的半径给定 $C$ 和 $D$.

\subsection{连续介质的作用量}
场 $u(\mathbf{x},t)$ 的作用量对时间和空间一起积分.时间两端的 $u$ 固定时,欧拉--拉格朗日方程把 $\partial_t$ 和 $\nabla$ 同等看待.

弦、膜或弹性体的动能减去势能,
\begin{equation}
  S=\int_{t_0}^{t_1}\int\Bigl(\half\rho u_t^2-\half K|\nabla u|^2+fu\Bigr)\,\mathrm{d}^n x\,\mathrm{d}t,
\end{equation}
给出
\begin{equation}
  \rho u_{tt}=\nabla\cdot(K\nabla u)+f.
\end{equation}
$\rho,K$ 是常数时这就是波动方程.对 $u_t$ 的动量密度是 $\pi=\rho u_t$,哈密顿量
\begin{equation}
  H=\int\Bigl(\half\rho u_t^2+\half K|\nabla u|^2-fu\Bigr)\,\mathrm{d}^n x
\end{equation}
在 $\rho,K,f$ 不显含 $t$ 时守恒.

矢量场若动能是 $\half\rho|\mathbf{u}_t|^2$,势能含 $|\nabla\otimes\mathbf{u}|^2$ 和 $(\nabla\cdot\mathbf{u})^2$,方程是弹性波.三维里 $|\nabla\otimes\mathbf{u}|^2=|\nabla\times\mathbf{u}|^2+|\nabla\cdot\mathbf{u}|^2$.取适当系数,作用量只保留 $|\mathbf{u}_t|^2-c^2|\nabla\times\mathbf{u}|^2$,得到无源麦克斯韦方程,$\mathbf{E}=\mathbf{u}_t$, $\mathbf{B}=\nabla\times\mathbf{u}$.

含二阶导数的作用量
\begin{equation}
  S=\int\int\Bigl(\half u_t^2-\half K(\Delta u)^2+fu\Bigr)\,\mathrm{d}^n x\,\mathrm{d}t
\end{equation}
给出 $u_{tt}+K\Delta^2 u=f$.$n=1$ 是梁,$n=2$ 是板.自由边上,一阶变分对 $\delta u$ 和 $\partial_\nu\delta u$ 都要为零,所以弯矩和等效剪力都是零.一维自由端就是 $u''=0$ 和 $u'''=0$.

复场 $\psi$ 的作用量
\begin{equation}
  S=\int_{t_0}^{t_1}\int\Biggl(
    \frac{\imath\hbar}{2}(\psi^*\psi_t-\psi\psi_t^*)
    -\frac{\hbar^2}{2m}|\nabla\psi|^2-V|\psi|^2
  \Biggr)\,\mathrm{d}^n x\,\mathrm{d}t
\end{equation}
对 $\psi^*$ 变分,得到
\begin{equation}
  \imath\hbar\psi_t=-\frac{\hbar^2}{2m}\Delta\psi+V\psi.
\end{equation}
这里 $\psi$ 和 $\psi^*$ 先当成独立的.哈密顿量是动能加势能的积分,不含 $\psi_t$,所以它就是能量.

平衡问题是作用量里没有时间的特例.狄利克雷积分的极小给出泊松方程,面积泛函的极小给出极小曲面,板的弯曲能 $\int(\Delta u)^2$ 的极小给出双调和方程 $\Delta^2 u=0$.边界上没固定的导数,一律由自然边界条件补上.

泛函的变分要写明哪些端点条件已经给定,哪些端点量仍可变化.乘子题要说明对乘子变分恢复的是哪一条约束.

\begin{problemset}
\renewcommand{\theenumi}{\thechapter.\arabic{enumi}}
\renewcommand{\labelenumi}{\theenumi.}
\item 对弧长
\begin{equation*}
  J[u]=\int_{x_1}^{x_2}\sqrt{1+(u')^2}\,\mathrm{d}x
\end{equation*}
写出欧拉--拉格朗日方程,并验证直线是解.端点 $(x_1,y_1)$, $(x_2,y_2)$ 固定.

\item 最速降线的时间是 $T[u]=\int\mathrm{d}s/v$.取适当坐标,使被积函数不显含自变量.写出贝尔特拉米第一积分,并说明解是摆线.

\item 设 $J[u]=\int_{x_1}^{x_2}F(x,u,u',u'')\,\mathrm{d}x$,端点的 $u$ 给定而 $u'$ 自由.导出自然边界条件,并说明它为什么是 $\partial F/\partial u''=0$.

\item 自由粒子 $L=\frac12 m\dot q^2$ 和谐振子 $L=\frac12 m\dot q^2-\frac12 m\omega^2 q^2$.写出哈密顿量,验证沿运动 $\mathrm{d}H/\mathrm{d}t=0$,并指出这个守恒对应 $L$ 的哪一种对称.

\item 用拉格朗日乘子求平面 $x+y+z=1$ 上离原点最近的点.

\item 周长固定的矩形,用乘子求面积的最大值,并核对它是正方形.

\item 悬链线:在长度固定的条件下,使重力势能取极值.写出带乘子的被积函数和欧拉--拉格朗日方程,说明解是双曲余弦.

\item 光线在平面上的速率只依赖高度,$c=c(u)$.旅行时间是 $\int_{x_1}^{x_2}\sqrt{1+(u')^2}/c(u)\,\mathrm{d}x$,两端固定.写出欧拉--拉格朗日方程,并用贝尔特拉米首次积分把它降成一阶方程.说明 $c$ 是常数时解是直线.

\item 膜的能量是面积减去外力做的功,
\begin{equation*}
  E=\iint_D\Bigl(k\sqrt{1+u_x^2+u_y^2}-fu\Bigr)\,\mathrm{d}x\,\mathrm{d}y,
\end{equation*}
边界上 $u$ 给定.写出欧拉--拉格朗日方程.$f=0$ 时它应当回到平均曲率为零.

\item 牛顿最小阻力问题里,被积函数是 $y(y')^3/(1+(y')^2)$,端点项是 $\half y(a)^2$.算出 $\partial L/\partial y'$,说明光滑自由端条件 $\partial L/\partial y'+y=0$ 没有实数解.再验证 $y'=-1$ 时 $\partial L/\partial y'=y$.

\item 区域 $0<x<2$, $0<y<1$.泛函
\begin{equation*}
  \Phi[u]=\iint\bigl(u_x^2+u_x u_y+u_y^2-2uf\bigr)\,\mathrm{d}x\,\mathrm{d}y,
\end{equation*}
$u$ 在 $x=0$ 和 $x=2$ 上为零,上下两边自由.写出方程,以及自由边上的自然边界条件.

\item 同一矩形上
\begin{equation*}
  \Phi[u]=\iint\bigl(u_x^2-u_x u_y+u_y^2\bigr)\,\mathrm{d}x\,\mathrm{d}y
  +\int_0^1 2ug(y)\,\mathrm{d}y,
\end{equation*}
线积分在 $x=0$ 上.$u$ 在 $y=0$, $y=1$ 和 $x=2$ 上为零.写出方程和 $x=0$ 上的边界条件.

\item $\Phi[u]=\int_0^1\bigl((u'')^2-2(u')^2+2f u\bigr)\,\mathrm{d}x+4u(0)^2$,已知 $u'(1)=u(1)$.写出欧拉--拉格朗日方程和 $x=0$ 处的自然边界条件.$x=1$ 的条件不要再改写成 $u''$ 的条件,它已经给定.

\item 梁的静弯曲能 $\Phi[u]=\int_0^{\ell}(u'')^2\,\mathrm{d}x$,两端的 $u$ 和 $u'$ 都自由.写出欧拉--拉格朗日方程,并由边界项说明自然边界条件是 $u''=u'''=0$.
\end{problemset}


% \section{变分法应用}
% \input{applications.tex}
